\documentclass[../../main.tex]{subfiles}% 注意这里的文件路径不能用 ./main.tex ,否则用latexmk编译子文件会报错
\graphicspath{{\subfix{./image/}}{\subfix{../../image/}}} % 指定图片引用路径,后续可以直接使用图片文件名
% 如果图片存放在上级目录,则使用 ../../image/ 作为备用路径

% 例如:
% \begin{figure}[H]
% \centering
% \includegraphics[width=0.3\textwidth]{图.png}
% \caption{}
% \label{figure:图}
% \end{figure}
% 注意:上述\label{}一定要放在\caption{}之后,否则引用图片序号会只会显示??.

\begin{document}
% SOURCE p=67; printed=52; object=ch02-complex-integral-title; region=bottom
\section{复值函数的积分}\label{section:ch02-complex-integration}
% SOURCE p=67,68; printed=52,53; object=ch02-complex-integral-intro; region=bottom/top
我们仍在一个固定的测度空间 $(X,\mathscr{M},\mu)$ 上讨论。上一节定义的积分可以用显然的方式推广到实值可测函数 $f$。
% SOURCE p=68; printed=53; object=ch02-real-integral; region=top
\begin{definition}[实值函数的积分]\label{definition:ch02-real-integral}
设 $f^+$ 和 $f^-$ 为 $f$ 的正部和负部。若 $\int f^+$ 和 $\int f^-$ 至少有一个有限，则定义
\begin{align*}
\int f=\int f^+-\int f^-.
\end{align*}
\end{definition}
% SOURCE p=68; printed=53; object=ch02-real-integrable; region=upper
\begin{definition}[实值可积函数]\label{definition:ch02-real-integrable}
我们主要关注 $\int f^+$ 和 $\int f^-$ 均有限的情形；此时称 $f$ \textbf{可积}。
\end{definition}
% SOURCE p=68; printed=53; object=ch02-real-integrable-absolute; region=upper
\begin{remark}
由于 $|f|=f^++f^-$，显然 $f$ 可积当且仅当 $\int|f|<\infty$。
\end{remark}
% SOURCE p=68; printed=53; object=ch02-real-integral-linearity; number=2.21; region=upper
\begin{proposition}\label{proposition:ch02-real-integral-linearity}
$X$ 上的实值可积函数构成一个实向量空间，积分是其上的线性泛函。
\end{proposition}
% SOURCE p=68; printed=53; object=ch02-proof-real-integral-linearity; region=upper-middle
\begin{proof}
第一个断言由 $|af+bg|\le|a||f|+|b||g|$ 得到。容易验证，对任意 $a\in\mathbb{R}$，$\int af=a\int f$。为证明可加性，设 $f$ 和 $g$ 可积，并令 $h=f+g$。则 $h^+-h^-=f^+-f^-+g^+-g^-$，因而 $h^++f^-+g^-=h^-+f^++g^+$。由\cref{theorem:ch02-integral-series}，
\begin{align*}
\int h^++\int f^-+\int g^-=\int h^-+\int f^++\int g^+.
\end{align*}
重新组合各项即得所需结论：
\begin{align*}
\int h=\int h^+-\int h^-=\int f^+-\int f^-+\int g^+-\int g^-=\int f+\int g.
\end{align*}
\end{proof}

% SOURCE p=68; printed=53; object=ch02-complex-integrable; region=middle
\begin{definition}[复值可积函数]\label{definition:ch02-complex-integrable}
若 $f$ 为复值可测函数，且 $\int|f|<\infty$，则称 $f$ \textbf{可积}。更一般地，若 $E\in\mathscr{M}$ 且 $\int_E|f|<\infty$，则称 $f$ \textbf{在 $E$ 上可积}。
\end{definition}
% SOURCE p=68; printed=53; object=ch02-complex-integrable-parts; region=middle
\begin{remark}
由于 $|f|\le|\operatorname{Re}f|+|\operatorname{Im}f|\le2|f|$，$f$ 可积当且仅当 $\operatorname{Re}f$ 和 $\operatorname{Im}f$ 均可积。
\end{remark}
% SOURCE p=68; printed=53; object=ch02-complex-integral; region=middle
\begin{definition}[复值函数的积分]\label{definition:ch02-complex-integral}
在 $f$ 可积时，定义
\begin{align*}
\int f=\int\operatorname{Re}f+\mathrm{i}\int\operatorname{Im}f.
\end{align*}
\end{definition}
% SOURCE p=68; printed=53; object=ch02-complex-integral-linearity; region=middle-lower
\begin{proposition}\label{proposition:ch02-complex-integral-linearity}
容易得出，复值可积函数构成一个复向量空间，积分是其上的复线性泛函。
\end{proposition}
% SOURCE p=68; printed=53; object=ch02-lone-provisional; region=lower
\begin{definition}[$L^1$ 的暂定记号]\label{definition:ch02-lone-provisional}
暂且将这个空间记为 $L^1(\mu)$（或依上下文记为 $L^1(X,\mu)$、$L^1(X)$，或简记为 $L^1$）。
\end{definition}
% SOURCE p=68; printed=53; object=ch02-lone-superscript; region=lower
上标 $1$ 是标准记号，它的含义要到\cref{chapter:ch06-lp-spaces}才会显现。
% SOURCE p=68; printed=53; object=ch02-integral-triangle; number=2.22; region=lower
\begin{proposition}\label{proposition:ch02-integral-triangle}
若 $f\in L^1$，则 $\left|\int f\right|\le\int|f|$。
\end{proposition}
% SOURCE p=68; printed=53; object=ch02-proof-integral-triangle; region=lower-bottom
\begin{proof}
若 $\int f=0$，结论显然。若 $f$ 为实值函数，结论也几乎显然，因为
\begin{align*}
\left|\int f\right|=\left|\int f^+-\int f^-\right|\le\int f^++\int f^-=\int|f|.
\end{align*}
若 $f$ 为复值函数且 $\int f\ne0$，令 $\alpha=\overline{\operatorname{sgn}(\int f)}$。则 $|\int f|=\alpha\int f=\int\alpha f$。特别地，$\int\alpha f$ 为实数，所以
\begin{align*}
\left|\int f\right|=\operatorname{Re}\int\alpha f=\int\operatorname{Re}(\alpha f)
\le\int|\operatorname{Re}(\alpha f)|\le\int|\alpha f|=\int|f|.
\end{align*}
\end{proof}

% SOURCE p=69; printed=54; object=ch02-lone-ae-equivalence; number=2.23; region=top
\begin{proposition}\label{proposition:ch02-lone-ae-equivalence}
\begin{enumerate}
\item\label{item:ch02-lone-sigma-finite-support} 若 $f\in L^1$，则 $\{x:f(x)\ne0\}$ 为 $\sigma$-有限集。
\item\label{item:ch02-lone-equality} 若 $f,g\in L^1$，则对所有 $E\in\mathscr{M}$ 都有 $\int_E f=\int_E g$，当且仅当 $\int|f-g|=0$，也当且仅当 $f=g$ 几乎处处成立。
\end{enumerate}
\end{proposition}
% SOURCE p=69; printed=54; object=ch02-proof-lone-ae-equivalence; region=upper
\begin{proof}
\subcref{proposition:ch02-lone-ae-equivalence}{item:ch02-lone-sigma-finite-support}以及\subcref{proposition:ch02-lone-ae-equivalence}{item:ch02-lone-equality}中的第二个等价关系，分别由\cref{proposition:ch02-finite-integral-support}和\cref{proposition:ch02-zero-integral}得到。若 $\int|f-g|=0$，则由\cref{proposition:ch02-integral-triangle}，对任意 $E\in\mathscr{M}$，
\begin{align*}
\left|\int_E f-\int_E g\right|\le\int\chi_E|f-g|\le\int|f-g|=0,
\end{align*}
因而 $\int_E f=\int_E g$。另一方面，若 $u=\operatorname{Re}(f-g)$、$v=\operatorname{Im}(f-g)$，且“$f=g$ 几乎处处成立”不成立，则 $u^+$、$u^-$、$v^+$ 和 $v^-$ 中至少有一个在某个正测度集上不为零。例如，若 $E=\{x:u^+(x)>0\}$ 具有正测度，则因 $E$ 上 $u^-=0$，有 $\operatorname{Re}(\int_E f-\int_E g)=\int_E u^+>0$；其他情形类似。
\end{proof}
% SOURCE p=69; printed=54; object=ch02-null-set-identification; region=middle
\begin{remark}
这一命题表明，就积分而言，在零测集上修改函数没有影响。事实上，对于仅在一个余集为零测集的可测集 $E$ 上有定义的函数 $f$，只需在 $E^c$ 上将 $f$ 定义为零（或任何其他值），便可对它积分。这样，在积分问题中，就可以将几乎处处有限的 $\overline{\mathbb{R}}$-值函数当作实值函数处理。
\end{remark}
% SOURCE p=69; printed=54; object=ch02-lone-equivalence-classes; region=middle
\begin{definition}[$L^1$ 空间]\label{definition:ch02-lone-equivalence-classes}
考虑到这一点，将 $L^1(\mu)$ 重新定义为 $X$ 上几乎处处有定义的可积函数的等价类组成的集合，会更加方便。其中 $f$ 和 $g$ 等价，当且仅当 $f=g$ 几乎处处成立。
\end{definition}
% SOURCE p=69; printed=54; object=ch02-lone-classes-notation; region=middle-lower
\begin{remark}
这个新的 $L^1(\mu)$ 仍是复向量空间（运算为几乎处处逐点的加法和数乘）。此后，虽然将 $L^1(\mu)$ 看作等价类的空间，仍用记号“$f\in L^1(\mu)$”表示 $f$ 是几乎处处有定义的可积函数。这种轻微的记号滥用是通常认可的，也很少引起混淆。
\end{remark}
% SOURCE p=69; printed=54; object=ch02-lone-completion; region=lower
\begin{remark}
$L^1(\mu)$ 的新定义还有两个优点。首先，若 $\overline{\mu}$ 为 $\mu$ 的完备化，\cref{proposition:ch02-measurable-version-completion}给出 $L^1(\overline{\mu})$ 与 $L^1(\mu)$ 之间自然的一一对应，所以可以（且我们将会）将这两个空间视为相同。
\end{remark}
% SOURCE p=69; printed=54; object=ch02-lone-metric; region=lower
\begin{proposition}\label{proposition:ch02-lone-metric}
其次，$L^1$ 是度量空间，其距离函数为 $\rho(f,g)=\int|f-g|$。
\end{proposition}
% SOURCE p=69; printed=54; object=ch02-proof-lone-metric; region=lower
\begin{proof}
三角不等式容易验证，且显然 $\rho(f,g)=\rho(g,f)$；但要使 $\rho(f,g)=0$ 仅在 $f=g$ 时成立，按照\subcref{proposition:ch02-lone-ae-equivalence}{item:ch02-lone-equality}，必须将几乎处处相等的函数视为相同。
\end{proof}
% SOURCE p=69; printed=54; object=ch02-lone-convergence; region=lower
\begin{definition}[$L^1$ 收敛]\label{definition:ch02-lone-convergence}
关于上述度量的收敛称为\textbf{$L^1$ 收敛}；因此，$f_n\to f$ 在 $L^1$ 中成立，当且仅当 $\int|f_n-f|\to0$。
\end{definition}
% SOURCE p=69; printed=54; object=ch02-dominated-convergence-intro; region=lower-bottom
下面给出三个基本收敛定理中的最后一个（另外两个是\cref{theorem:ch02-monotone-convergence}和\cref{lemma:ch02-fatou}），并推出一些有用的结论。在 Fatou 引理之前讨论的 $\mathbb{R}$ 上 Lebesgue 测度积分的背景下，这一定理的思想是：若 $f_n\to f$ 几乎处处成立，且 $|f_n|$ 的图形被限制在平面中一个面积有限的区域内，使得图形下方的面积不能逃向无穷远，则 $\int f_n\to\int f$。
% SOURCE p=69; printed=54; object=ch02-dominated-convergence; number=2.24; region=bottom
\begin{theorem}[控制收敛定理]\label{theorem:ch02-dominated-convergence}
设 $\{f_n\}$ 为 $L^1$ 中的序列，满足
\begin{enumerate}
\item\label{item:ch02-dominated-pointwise} $f_n\to f$ 几乎处处成立；
\item\label{item:ch02-dominated-majorant} 存在非负函数 $g\in L^1$，使得对所有 $n$，$|f_n|\le g$ 几乎处处成立。
\end{enumerate}
则 $f\in L^1$，且 $\int f=\lim_{n\to\infty}\int f_n$。
\end{theorem}
% SOURCE p=70; printed=55; object=ch02-proof-dominated-convergence; region=top
\begin{proof}
由\cref{proposition:ch02-complete-ae-measurability,proposition:ch02-measurable-version-completion}，$f$ 可测（可能需要在一个零测集上重新定义）。又因 $|f|\le g$ 几乎处处成立，故 $f\in L^1$。分别取实部与虚部，只需假定 $f_n$ 和 $f$ 为实值函数；此时 $g+f_n\ge0$ 及 $g-f_n\ge0$ 都几乎处处成立。因此由\cref{lemma:ch02-fatou}，
\begin{align*}
\int g+\int f&\le\liminf\int(g+f_n)=\int g+\liminf\int f_n,\\
\int g-\int f&\le\liminf\int(g-f_n)=\int g-\limsup\int f_n.
\end{align*}
所以 $\liminf\int f_n\ge\int f\ge\limsup\int f_n$，结论随之成立。
\end{proof}
% SOURCE p=70; printed=55; object=ch02-absolutely-integrable-series; number=2.25; region=upper-middle
\begin{theorem}\label{theorem:ch02-absolutely-integrable-series}
设 $\{f_j\}$ 为 $L^1$ 中的序列，且 $\sum_1^\infty\int|f_j|<\infty$。则 $\sum_1^\infty f_j$ 几乎处处收敛于一个 $L^1$ 中的函数，且 $\int\sum_1^\infty f_j=\sum_1^\infty\int f_j$。
\end{theorem}
% SOURCE p=70; printed=55; object=ch02-proof-absolutely-integrable-series; region=middle
\begin{proof}
由\cref{theorem:ch02-integral-series}，$\int\sum_1^\infty|f_j|=\sum_1^\infty\int|f_j|<\infty$，所以函数 $g=\sum_1^\infty|f_j|$ 属于 $L^1$。特别地，由\cref{proposition:ch02-finite-integral-support}，对几乎处处的 $x$，$\sum_1^\infty|f_j(x)|$ 有限，对每个这样的 $x$，级数 $\sum_1^\infty f_j(x)$ 收敛。此外，对所有 $n$，$|\sum_1^n f_j|\le g$，故可对部分和序列应用\cref{theorem:ch02-dominated-convergence}，得到 $\int\sum_1^\infty f_j=\sum_1^\infty\int f_j$。
\end{proof}
% SOURCE p=70; printed=55; object=ch02-lone-density; number=2.26; region=middle
\begin{theorem}\label{theorem:ch02-lone-density}
若 $f\in L^1(\mu)$ 且 $\epsilon>0$，则存在可积简单函数 $\phi=\sum a_j\chi_{E_j}$，使 $\int|f-\phi|\,\mathrm{d}\mu<\epsilon$。（也就是说，可积简单函数在 $L^1$ 度量下稠密于 $L^1$。）若 $\mu$ 为 $\mathbb{R}$ 上的 Lebesgue--Stieltjes 测度，则 $\phi$ 定义中的集合 $E_j$ 可取为有限个开区间的并；此外，存在一个在某个有界区间之外为零的连续函数 $g$，使 $\int|f-g|\,\mathrm{d}\mu<\epsilon$。
\end{theorem}
% SOURCE p=70; printed=55; object=ch02-proof-lone-density; region=lower
\begin{proof}
取\subcref{theorem:ch02-simple-approximation}{item:ch02-simple-approximation-complex}中的 $\{\phi_n\}$。由于 $|\phi_n-f|\le2|f|$，由\cref{theorem:ch02-dominated-convergence}，当 $n$ 充分大时，$\int|\phi_n-f|<\epsilon$。若 $\phi_n=\sum a_j\chi_{E_j}$，其中 $E_j$ 两两不交且 $a_j$ 非零，则
\begin{align*}
\mu(E_j)=|a_j|^{-1}\int_{E_j}|\phi_n|\le|a_j|^{-1}\int|f|<\infty.
\end{align*}
此外，若 $E$ 和 $F$ 为可测集，则 $\mu(E\mathbin{\triangle}F)=\int|\chi_E-\chi_F|$。所以，若 $\mu$ 为 $\mathbb{R}$ 上的 Lebesgue--Stieltjes 测度，由\cref{proposition:ch01-finite-interval-approximation}，可以在 $L^1$ 度量下用函数 $\chi_{I_k}$ 的有限和任意精确地逼近 $\chi_{E_j}$，其中 $I_k$ 为开区间。最后，若 $I_k=(a,b)$，则可用在 $(a,b)$ 之外为零的连续函数在 $L^1$ 度量下逼近 $\chi_{I_k}$。
\begin{example*}
给定 $\epsilon>0$，取连续函数 $g$，使其在 $(-\infty,a]$ 和 $[b,\infty)$ 上等于 $0$，在 $[a+\epsilon,b-\epsilon]$ 上等于 $1$，在 $[a,a+\epsilon]$ 和 $[b-\epsilon,b]$ 上为线性函数。
\end{example*}
综合这些事实，即得所需断言。
\end{proof}
% SOURCE p=70; printed=55; object=ch02-parameter-integrals-intro; region=bottom
下一个定理给出交换极限或导数与积分的一个判别准则，其限制比大多数高等微积分教材中的条件更弱。

% SOURCE p=71; printed=56; object=ch02-parameter-integrals; number=2.27; region=top
\begin{theorem}\label{theorem:ch02-parameter-integrals}
设 $f:X\times[a,b]\to\mathbb{C}$，其中 $-\infty<a<b<\infty$，且对每个 $t\in[a,b]$，$f(\cdot,t):X\to\mathbb{C}$ 可积。令 $F(t)=\int_X f(x,t)\,\mathrm{d}\mu(x)$。
\begin{enumerate}
\item\label{item:ch02-parameter-continuity} 假设存在 $g\in L^1(\mu)$，使得对所有 $x,t$，$|f(x,t)|\le g(x)$。若对每个 $x$，$\lim_{t\to t_0}f(x,t)=f(x,t_0)$，则 $\lim_{t\to t_0}F(t)=F(t_0)$；特别地，若对每个 $x$，$f(x,\cdot)$ 连续，则 $F$ 连续。
\item\label{item:ch02-parameter-differentiation} 假设 $\frac{\partial f}{\partial t}$ 存在，且存在 $g\in L^1(\mu)$，使得对所有 $x,t$，$|\frac{\partial f}{\partial t}(x,t)|\le g(x)$。则 $F$ 可微，且 $F'(t)=\int\frac{\partial f}{\partial t}(x,t)\,\mathrm{d}\mu(x)$。
\end{enumerate}
\end{theorem}
% SOURCE p=71; printed=56; object=ch02-proof-parameter-integrals; region=upper-middle
\begin{proof}
对于\subcref{theorem:ch02-parameter-integrals}{item:ch02-parameter-continuity}，对 $f_n(x)=f(x,t_n)$ 应用\cref{theorem:ch02-dominated-convergence}，其中 $\{t_n\}$ 为 $[a,b]$ 中任意收敛于 $t_0$ 的序列。对于\subcref{theorem:ch02-parameter-integrals}{item:ch02-parameter-differentiation}，注意到
\begin{align*}
\frac{\partial f}{\partial t}(x,t_0)=\lim h_n(x),\qquad
h_n(x)=\frac{f(x,t_n)-f(x,t_0)}{t_n-t_0},
\end{align*}
其中 $\{t_n\}$ 仍为任意收敛于 $t_0$ 的序列。因此 $\frac{\partial f}{\partial t}$ 可测，且由中值定理，
\begin{align*}
|h_n(x)|\le\sup_{t\in[a,b]}\left|\frac{\partial f}{\partial t}(x,t)\right|\le g(x).
\end{align*}
于是可以再次应用\cref{theorem:ch02-dominated-convergence}，得到
\begin{align*}
F'(t_0)=\lim\frac{F(t_n)-F(t_0)}{t_n-t_0}=\lim\int h_n(x)\,\mathrm{d}\mu(x)=\int\frac{\partial f}{\partial t}(x,t)\,\mathrm{d}\mu(x).
\end{align*}
\end{proof}
% SOURCE p=71; printed=56; object=ch02-sequential-limit-remark; region=middle
\begin{remark}
上述证明使用收敛于 $t_0$ 的序列，在技术上是必要的，因为控制收敛定理只处理函数序列。不过，在这种情形下，我们通常只说“令 $t\to t_0$”，并约定论证的基础是序列收敛。
\end{remark}
% SOURCE p=71; printed=56; object=ch02-local-parameter-remark; region=middle-lower
\begin{remark}
需要注意，\cref{theorem:ch02-parameter-integrals}中 $f$ 或 $\frac{\partial f}{\partial t}$ 的估计成立的区间 $[a,b]$，可能只是 $f(x,\cdot)$ 的定义域——某个开区间 $I$（也可能是 $\mathbb{R}$ 本身）——的真子区间。若\subcref{theorem:ch02-parameter-integrals}{item:ch02-parameter-continuity}或\subcref{theorem:ch02-parameter-integrals}{item:ch02-parameter-differentiation}的假设对所有 $[a,b]\subseteq I$ 成立，控制函数 $g$ 可以依赖于 $a,b$，则积分所得函数 $F$ 在整个 $I$ 上连续或可微，因为这些性质本质上是局部的。
\end{remark}
% SOURCE p=71; printed=56; object=ch02-lebesgue-integral-name; region=lower
\begin{definition}[Lebesgue 积分]\label{definition:ch02-lebesgue-integral}
当测度 $\mu$ 为 $\mathbb{R}$ 上的 Lebesgue 测度时，我们所建立的积分称为 \textbf{Lebesgue 积分}。
\end{definition}
% SOURCE p=71; printed=56; object=ch02-riemann-comparison-intro; region=lower
现在适合研究 $\mathbb{R}$ 上 Lebesgue 积分与 Riemann 积分的关系。我们将使用 Darboux 通过上、下和对 Riemann 积分给出的刻画，现回顾如下。
% SOURCE p=71; printed=56; object=ch02-interval-partition; region=lower
\begin{definition}[区间的分割]\label{definition:ch02-interval-partition}
设 $[a,b]$ 为紧区间。$[a,b]$ 的一个\textbf{分割}是有限序列 $P=\{t_j\}_0^n$，满足 $a=t_0<t_1<\cdots<t_n=b$。
\end{definition}
% SOURCE p=71,72; printed=56,57; object=ch02-darboux-sums; region=bottom/top
\begin{definition}[上和与下和]\label{definition:ch02-darboux-sums}
设 $f$ 为 $[a,b]$ 上任意有界实值函数。对每个分割 $P$，定义
\begin{align*}
S_Pf=\sum_1^n M_j(t_j-t_{j-1}),\qquad
s_Pf=\sum_1^n m_j(t_j-t_{j-1}),
\end{align*}
其中 $M_j$ 和 $m_j$ 分别为 $f$ 在 $[t_{j-1},t_j]$ 上的上确界和下确界。
\end{definition}
% SOURCE p=72; printed=57; object=ch02-riemann-integral; region=top
\begin{definition}[Riemann 积分]\label{definition:ch02-riemann-integral}
定义
\begin{align*}
\overline{I}_a^b(f)=\inf_P S_Pf,\qquad
\underline{I}_a^b(f)=\sup_P s_Pf,
\end{align*}
其中下确界和上确界遍历所有分割 $P$。若 $\overline{I}_a^b(f)=\underline{I}_a^b(f)$，则它们的公共值就是 \textbf{Riemann 积分} $\int_a^b f(x)\,\mathrm{d}x$，并称 $f$ \textbf{Riemann 可积}。
\end{definition}
% SOURCE p=72; printed=57; object=ch02-riemann-lebesgue; number=2.28; region=upper
\begin{theorem}\label{theorem:ch02-riemann-lebesgue}
设 $f$ 为 $[a,b]$ 上的有界实值函数。
\begin{enumerate}
\item\label{item:ch02-riemann-lebesgue-equality} 若 $f$ Riemann 可积，则 $f$ Lebesgue 可测（又因它有界，故在 $[a,b]$ 上可积），且 $\int_a^b f(x)\,\mathrm{d}x=\int_{[a,b]}f\,\mathrm{d}m$。
\item\label{item:ch02-riemann-discontinuities} $f$ Riemann 可积，当且仅当 $\{x\in[a,b]:f\text{ 在 }x\text{ 处不连续}\}$ 的 Lebesgue 测度为零。
\end{enumerate}
\end{theorem}
% SOURCE p=72; printed=57; object=ch02-proof-riemann-lebesgue; region=middle
\begin{proof}
假设 $f$ Riemann 可积。对每个分割 $P$，令
\begin{align*}
G_P=\sum_1^n M_j\chi_{(t_{j-1},t_j]},\qquad
g_P=\sum_1^n m_j\chi_{(t_{j-1},t_j]},
\end{align*}
记号同前。于是 $S_Pf=\int G_P\,\mathrm{d}m$ 且 $s_Pf=\int g_P\,\mathrm{d}m$。
\begin{definition*}[分割的网径]
分割的\textbf{网径}是 $\max_j(t_j-t_{j-1})$。
\end{definition*}
存在一个网径趋于零的分割序列 $\{P_k\}$，每个分割包含前一个分割（所以 $g_{P_k}$ 随 $k$ 递增，而 $G_{P_k}$ 递减），并使 $S_{P_k}f$ 和 $s_{P_k}f$ 均收敛于 $\int_a^b f(x)\,\mathrm{d}x$。令 $G=\lim G_{P_k}$ 和 $g=\lim g_{P_k}$。则 $g\le f\le G$，由\cref{theorem:ch02-dominated-convergence}，
\begin{align*}
\int G\,\mathrm{d}m=\int g\,\mathrm{d}m=\int_a^b f(x)\,\mathrm{d}x.
\end{align*}
因而 $\int(G-g)\,\mathrm{d}m=0$，由\cref{proposition:ch02-zero-integral}，$G=g$ 几乎处处成立，从而 $G=f$ 几乎处处成立。$G$ 是简单函数序列的极限，故可测；又因 $m$ 完备，所以 $f$ 可测，且
\begin{align*}
\int_{[a,b]}f\,\mathrm{d}m=\int G\,\mathrm{d}m=\int_a^b f(x)\,\mathrm{d}x.
\end{align*}
这证明了\subcref{theorem:ch02-riemann-lebesgue}{item:ch02-riemann-lebesgue-equality}，而\subcref{theorem:ch02-riemann-lebesgue}{item:ch02-riemann-discontinuities}的证明提纲见\cref{exercise:ch02-riemann-discontinuities}。
\end{proof}
% SOURCE p=72; printed=57; object=ch02-improper-integrals-discussion; region=lower
由此，（正常的）Riemann 积分包含在 Lebesgue 积分之中。某些广义 Riemann 积分（绝对收敛的那些）可以直接解释为 Lebesgue 积分，另一些则仍然需要取极限。
% SOURCE p=72; printed=57; object=ch02-improper-integral-example; region=lower
\begin{example}\label{example:ch02-improper-integral}
若对所有 $b>0$，$f$ 在 $[0,b]$ 上 Riemann 可积，且在 $[0,\infty)$ 上 Lebesgue 可积，则由\cref{theorem:ch02-dominated-convergence}，
\begin{align*}
\int_{[0,\infty)}f\,\mathrm{d}m=\lim_{b\to\infty}\int_0^b f(x)\,\mathrm{d}x.
\end{align*}
但即使 $f$ 不可积，右端极限也可能存在。例如，$f=\sum_1^\infty n^{-1}(-1)^n\chi_{(n,n+1]}$。
\end{example}
% SOURCE p=72; printed=57; object=ch02-lebesgue-interval-notation; region=lower
此后，我们通常将记号 $\int_a^b f(x)\,\mathrm{d}x$ 用于 Lebesgue 积分。
% SOURCE p=72,73; printed=57,58; object=ch02-integral-constructions-comparison; region=bottom/top
\begin{remark}
比较 Lebesgue 积分与 Riemann 积分的构造，或许有助于理解。设 $f$ 为 $[a,b]$ 上的有界可测函数，为简单起见，假设 $f\ge0$。计算 $f$ 的 Riemann 积分时，将区间 $[a,b]$ 分割为子区间，并用在每个子区间上为常数的函数从上、下逼近 $f$。计算 $f$ 的 Lebesgue 积分时，则选取递增趋于 $f$ 的简单函数序列。特别地，若选取\subcref{theorem:ch02-simple-approximation}{item:ch02-simple-approximation-positive}证明中构造的序列（见\cref{figure:ch02-simple-approximation}），实质上是在将 $f$ 的值域分割为子区间 $I_j$，并在每个集合 $f^{-1}(I_j)$ 上用常数逼近 $f$。由于即使 $f$ 连续，集合 $f^{-1}(I_j)$ 也可能很复杂，这种方法一开始便需要更精细的测度理论；但它更能适应所考虑的具体函数 $f$，因此更加灵活，也更容易推广。（在 Lebesgue 理论中，假定 $f$ 可测，使得不必同时考虑上、下两种逼近；不过，后一种观点在抽象背景下也可以采用，见\cref{exercise:ch02-upper-lower-simple}。）
\end{remark}
% SOURCE p=73; printed=58; object=ch02-lebesgue-advantages; region=upper-middle
Lebesgue 理论相较于 Riemann 理论有两个实质优点。首先，有更强的收敛定理可用，例如单调收敛定理和控制收敛定理。这些定理不仅给出以前无法得到的结果，也减少了证明经典定理所需的工作。其次，可以积分的函数范围更广。
% SOURCE p=73; printed=58; object=ch02-rational-indicator-example; region=middle
\begin{example}\label{example:ch02-rational-indicator}
若 $R$ 为 $[0,1]$ 中有理数组成的集合，则 $\chi_R$ 在 $[0,1]$ 上处处不连续，故不 Riemann 可积；但它 Lebesgue 可积，且 $\int\chi_R\,\mathrm{d}m=0$。（实际上，这在某种意义上是个平凡例子，因为 $\chi_R$ 几乎处处等于常值函数 $0$。更有意思的例子见\cref{exercise:ch02-dense-singularities}。）
\end{example}
% SOURCE p=73; printed=58; object=ch02-completeness-benefit; region=middle
当然，在经典分析中遇到的几乎所有函数都（局部）Riemann 可积，所以计算具体积分时，很少用到这种更大的普遍性。然而，它有一个关键作用：对于由积分定义度量的各种函数度量空间，使用 Lebesgue 可积函数时空间是完备的，而仅考虑 Riemann 可积函数时则不完备。以后会更深入地研究这一情形，尤其是在\cref{chapter:ch06-lp-spaces}中。（我们实际上已经通过\cref{theorem:ch02-absolutely-integrable-series}的一种隐蔽形式证明了 $L^1(\mu)$ 的完备性。要揭示这种形式，参见\cref{theorem:ch05-banach-series-criterion}。）
% SOURCE p=73; printed=58; object=ch02-gamma-intro; region=middle-lower

本节最后介绍高等超越函数中最常见的一个，即 \textbf{gamma 函数} $\Gamma$；它将在后面的多处发挥作用。
% SOURCE p=73; printed=58; object=ch02-gamma-integrability; region=lower

若 $z\in\mathbb{C}$ 且 $\operatorname{Re}z>0$，定义 $f_z:(0,\infty)\to\mathbb{C}$ 为 $f_z(t)=t^{z-1}e^{-t}$。（这里 $t^{z-1}=\exp[(z-1)\log t]$。）由于 $|t^{z-1}|=t^{\operatorname{Re}z-1}$，有 $|f_z(t)|\le t^{\operatorname{Re}z-1}$；且在 $t\ge1$ 时，$|f_z(t)|\le C_ze^{-t/2}$。将 $t^{\operatorname{Re}z-1}e^{-t/2}$ 最大化，容易求出 $C_z$ 的精确值，不过此处并不需要。由于当 $a>-1$ 时，$\int_0^1t^a\,\mathrm{d}t<\infty$，且 $\int_1^\infty e^{-t/2}\,\mathrm{d}t<\infty$，可见当 $\operatorname{Re}z>0$ 时，$f_z\in L^1((0,\infty))$。
% SOURCE p=73; printed=58; object=ch02-gamma-definition; region=lower
\begin{definition}[gamma 函数]\label{definition:ch02-gamma}
定义
\begin{align*}
\Gamma(z)=\int_0^\infty t^{z-1}e^{-t}\,\mathrm{d}t\qquad(\operatorname{Re}z>0).
\end{align*}
\end{definition}
% SOURCE p=73; printed=58; object=ch02-gamma-functional-equation; region=lower-bottom
\begin{proposition}\label{proposition:ch02-gamma-functional-equation}
当 $\operatorname{Re}z>0$ 时，$\Gamma$ 满足函数方程
\begin{align*}
\Gamma(z+1)=z\Gamma(z).
\end{align*}
\end{proposition}
\begin{proof}
由分部积分，
\begin{align*}
\int_\epsilon^N t^ze^{-t}\,\mathrm{d}t=-t^ze^{-t}\Big|_\epsilon^N+z\int_\epsilon^N t^{z-1}e^{-t}\,\mathrm{d}t.
\end{align*}
令 $\epsilon\to0$ 且 $N\to\infty$，即可得到上述函数方程。
\end{proof}
% SOURCE p=73,74; printed=58,59; object=ch02-gamma-extension; region=bottom/top
\begin{definition}[gamma 函数的延拓]\label{definition:ch02-gamma-extension}
利用这个方程，可以将 $\Gamma$ 延拓到（几乎）整个复平面。具体地，当 $-1<\operatorname{Re}z\le0$ 时，定义 $\Gamma(z)=\frac{\Gamma(z+1)}{z}$；再用归纳法，若已在 $\operatorname{Re}z>-n$ 上定义 $\Gamma(z)$，就在 $\operatorname{Re}z>-n-1$ 上定义 $\Gamma(z)=\frac{\Gamma(z+1)}{z}$。这样得到的函数在整个 $\mathbb{C}$ 上有定义，但非正整数处是奇点，因为上述算法在这些点涉及除以零。
\end{definition}
% SOURCE p=74; printed=59; object=ch02-gamma-factorial; region=top
\begin{proposition}\label{proposition:ch02-gamma-factorial}
$\Gamma(1)=1$，且 $\Gamma(n+1)=n!$。
\end{proposition}
% SOURCE p=74; printed=59; object=ch02-proof-gamma-factorial; region=top
\begin{proof}
有 $\Gamma(1)=\int_0^\infty e^{-t}\,\mathrm{d}t=-e^{-t}|_0^\infty=1$；将函数方程应用 $n$ 次，即得 $\Gamma(n+1)=n!$。
\end{proof}
% SOURCE p=74; printed=59; object=ch02-gamma-factorial-remark; region=upper
（这一事实的另一种证明提纲见\cref{exercise:ch02-gamma-gaussian-moments}。）gamma 函数的许多应用，都用到它将阶乘函数推广到非整数这一事实。

\subsection*{习题}
% SOURCE p=74; printed=59; object=ch02-exercise-fatou-lower-bound; number=18; region=upper
\begin{exercise}\label{exercise:ch02-fatou-lower-bound}
若将假设 $f_n\in L^+$ 改为 $f_n$ 可测且 $f_n\ge-g$，其中 $g\in L^+\cap L^1$，则\cref{lemma:ch02-fatou}仍成立。对于非正函数，Fatou 引理的相应形式是什么？
\end{exercise}
% SOURCE p=74; printed=59; object=ch02-exercise-uniform-integral-limit; number=19; region=upper-middle
\begin{exercise}\label{exercise:ch02-uniform-integral-limit}
设 $\{f_n\}\subseteq L^1(\mu)$，且 $f_n\to f$ 一致收敛。
\begin{enumerate}
\item\label{item:ch02-uniform-finite-measure} 若 $\mu(X)<\infty$，则 $f\in L^1(\mu)$，且 $\int f_n\to\int f$。
\item\label{item:ch02-uniform-infinite-measure} 若 $\mu(X)=\infty$，则\subcref{exercise:ch02-uniform-integral-limit}{item:ch02-uniform-finite-measure}的结论可能不成立。（在带 Lebesgue 测度的 $\mathbb{R}$ 上找出例子。）
\end{enumerate}
\end{exercise}
% SOURCE p=74; printed=59; object=ch02-exercise-generalized-dominated; number=20; region=middle
\begin{exercise}[广义控制收敛定理]\label{exercise:ch02-generalized-dominated}
若 $f_n,g_n,f,g\in L^1$，$f_n\to f$ 且 $g_n\to g$ 几乎处处成立，$|f_n|\le g_n$，并且 $\int g_n\to\int g$，则 $\int f_n\to\int f$。（修改控制收敛定理的证明。）
\end{exercise}
% SOURCE p=74; printed=59; object=ch02-exercise-scheffe; number=21; region=middle
\begin{exercise}\label{exercise:ch02-scheffe}
设 $f_n,f\in L^1$，且 $f_n\to f$ 几乎处处成立。则 $\int|f_n-f|\to0$ 当且仅当 $\int|f_n|\to\int|f|$。（使用\cref{exercise:ch02-generalized-dominated}。）
\end{exercise}
% SOURCE p=74; printed=59; object=ch02-exercise-counting-convergence; number=22; region=middle
\begin{exercise}\label{exercise:ch02-counting-convergence}
设 $\mu$ 为 $\mathbb{N}$ 上的计数测度。将\cref{lemma:ch02-fatou}、\cref{theorem:ch02-monotone-convergence}和\cref{theorem:ch02-dominated-convergence}解释为关于无穷级数的命题。
\end{exercise}
% SOURCE p=74; printed=59; object=ch02-exercise-riemann-discontinuities; number=23; region=middle-lower
\begin{exercise}\label{exercise:ch02-riemann-discontinuities}
给定有界函数 $f:[a,b]\to\mathbb{R}$，令
\begin{align*}
H(x)=\lim_{\delta\to0}\sup_{|y-x|\le\delta}f(y),\qquad
h(x)=\lim_{\delta\to0}\inf_{|y-x|\le\delta}f(y).
\end{align*}
通过建立下面的引理，证明\subcref{theorem:ch02-riemann-lebesgue}{item:ch02-riemann-discontinuities}：
\begin{enumerate}
\item\label{item:ch02-oscillation-continuity} $H(x)=h(x)$ 当且仅当 $f$ 在 $x$ 处连续。
\item\label{item:ch02-envelope-integrals} 采用\subcref{theorem:ch02-riemann-lebesgue}{item:ch02-riemann-lebesgue-equality}证明中的记号，$H=G$ 和 $h=g$ 都几乎处处成立。因此，$H$ 和 $h$ 为 Lebesgue 可测的，且 $\int_{[a,b]}H\,\mathrm{d}m=\overline{I}_a^b(f)$，$\int_{[a,b]}h\,\mathrm{d}m=\underline{I}_a^b(f)$。
\end{enumerate}
\end{exercise}
% SOURCE p=74; printed=59; object=ch02-exercise-upper-lower-simple; number=24; region=lower
\begin{exercise}\label{exercise:ch02-upper-lower-simple}
设 $(X,\mathscr{M},\mu)$ 为测度空间，$\mu(X)<\infty$，其完备化为 $(X,\overline{\mathscr{M}},\overline{\mu})$。设 $f:X\to\mathbb{R}$ 有界。则 $f$ 为 $\overline{\mathscr{M}}$-可测的（因而属于 $L^1(\overline{\mu})$），当且仅当存在 $\mathscr{M}$-可测简单函数序列 $\{\phi_n\}$ 和 $\{\psi_n\}$，使 $\phi_n\le f\le\psi_n$，且 $\int(\psi_n-\phi_n)\,\mathrm{d}\mu<n^{-1}$。此时，$\lim\int\phi_n\,\mathrm{d}\mu=\lim\int\psi_n\,\mathrm{d}\mu=\int f\,\mathrm{d}\overline{\mu}$。
\end{exercise}
% SOURCE p=74,75; printed=59,60; object=ch02-exercise-dense-singularities; number=25; region=bottom/top
\begin{exercise}\label{exercise:ch02-dense-singularities}
当 $0<x<1$ 时，令 $f(x)=x^{-1/2}$，其余情形令 $f(x)=0$。设 $\{r_n\}_1^\infty$ 为有理数的一个枚举，并令 $g(x)=\sum_1^\infty2^{-n}f(x-r_n)$。
\begin{enumerate}
\item\label{item:ch02-dense-singularities-integrable} $g\in L^1(m)$；特别地，$g<\infty$ 几乎处处成立。
\item\label{item:ch02-dense-singularities-unbounded} $g$ 处处不连续，且在每个区间上都无界；在任意 Lebesgue 零测集上修改它之后，仍然如此。
\item\label{item:ch02-dense-singularities-square} $g^2<\infty$ 几乎处处成立，但 $g^2$ 在任何区间上都不可积。
\end{enumerate}
\end{exercise}
% SOURCE p=75; printed=60; object=ch02-exercise-indefinite-continuity; number=26; region=upper
\begin{exercise}\label{exercise:ch02-indefinite-continuity}
若 $f\in L^1(m)$，且 $F(x)=\int_{-\infty}^x f(t)\,\mathrm{d}t$，则 $F$ 在 $\mathbb{R}$ 上连续。
\end{exercise}
% SOURCE p=75; printed=60; object=ch02-exercise-series-interchange-failure; number=27; region=upper
\begin{exercise}\label{exercise:ch02-series-interchange-failure}
令 $f_n(x)=ae^{-nax}-be^{-nbx}$，其中 $0<a<b$。
\begin{enumerate}
\item\label{item:ch02-series-absolute-divergence} $\sum_1^\infty\int_0^\infty|f_n(x)|\,\mathrm{d}x=\infty$。
\item\label{item:ch02-sum-integrals-zero} $\sum_1^\infty\int_0^\infty f_n(x)\,\mathrm{d}x=0$。
\item\label{item:ch02-integral-sum-log} $\sum_1^\infty f_n\in L^1([0,\infty),m)$，且 $\int_0^\infty\sum_1^\infty f_n(x)\,\mathrm{d}x=\log\frac{b}{a}$。
\end{enumerate}
\end{exercise}
% SOURCE p=75; printed=60; object=ch02-exercise-compute-limits; number=28; region=upper-middle
\begin{exercise}\label{exercise:ch02-compute-limits}
计算下列极限，并说明计算的依据：
\begin{enumerate}
\item\label{item:ch02-limit-exponential-sine} $\displaystyle\lim_{n\to\infty}\int_0^\infty\left(1+\frac{x}{n}\right)^{-n}\sin\frac{x}{n}\,\mathrm{d}x$。
\item\label{item:ch02-limit-rational-powers} $\displaystyle\lim_{n\to\infty}\int_0^1(1+nx^2)(1+x^2)^{-n}\,\mathrm{d}x$。
\item\label{item:ch02-limit-sine-rational} $\displaystyle\lim_{n\to\infty}\int_0^\infty n\sin\frac{x}{n}\,[x(1+x^2)]^{-1}\,\mathrm{d}x$。
\item\label{item:ch02-limit-cauchy-kernel} $\displaystyle\lim_{n\to\infty}\int_a^\infty n(1+n^2x^2)^{-1}\,\mathrm{d}x$。（答案取决于 $a>0$、$a=0$ 还是 $a<0$。这与各种收敛定理如何相符？）
\end{enumerate}
\end{exercise}
% SOURCE p=75; printed=60; object=ch02-exercise-gamma-gaussian-moments; number=29; region=middle
\begin{exercise}\label{exercise:ch02-gamma-gaussian-moments}
通过对等式 $\int_0^\infty e^{-tx}\,\mathrm{d}x=\frac1t$ 求导，证明 $\int_0^\infty x^ne^{-x}\,\mathrm{d}x=n!$。类似地，通过对等式 $\int_{-\infty}^\infty e^{-tx^2}\,\mathrm{d}x=\sqrt{\frac\pi t}$ 求导，证明 $\int_{-\infty}^\infty x^{2n}e^{-x^2}\,\mathrm{d}x=\frac{(2n)!\sqrt\pi}{4^n n!}$（见\cref{proposition:ch02-gaussian-integral}）。
\end{exercise}
% SOURCE p=75; printed=60; object=ch02-exercise-factorial-limit; number=30; region=middle
\begin{exercise}\label{exercise:ch02-factorial-limit}
证明 $\displaystyle\lim_{k\to\infty}\int_0^k x^n(1-k^{-1}x)^k\,\mathrm{d}x=n!$。
\end{exercise}
% SOURCE p=75; printed=60; object=ch02-exercise-termwise-integrals; number=31; region=lower
\begin{exercise}\label{exercise:ch02-termwise-integrals}
将被积函数的一部分展开为无穷级数，并论证逐项积分的合理性，由此推导下列公式。\cref{exercise:ch02-gamma-gaussian-moments}可能有用。（注意：在\subcref{exercise:ch02-termwise-integrals}{item:ch02-laplace-sine}和\subcref{exercise:ch02-termwise-integrals}{item:ch02-laplace-bessel}中，只有 $a>1$ 时才能逐项积分，且所得级数收敛；但所列公式实际上对所有 $a>0$ 都成立。）
\begin{enumerate}
\item\label{item:ch02-gaussian-cosine} 对 $a>0$，$\int_{-\infty}^\infty e^{-x^2}\cos ax\,\mathrm{d}x=\sqrt\pi e^{-a^2/4}$。
\item\label{item:ch02-log-power-integral} 对 $a>-1$，$\int_0^1x^a(1-x)^{-1}\log x\,\mathrm{d}x=-\sum_1^\infty(a+k)^{-2}$。
\item\label{item:ch02-gamma-zeta-integral} 对 $a>1$，$\int_0^\infty x^{a-1}(e^x-1)^{-1}\,\mathrm{d}x=\Gamma(a)\zeta(a)$，其中 $\zeta(a)=\sum_1^\infty n^{-a}$。
\item\label{item:ch02-laplace-sine} 对 $a>1$，$\int_0^\infty e^{-ax}x^{-1}\sin x\,\mathrm{d}x=\arctan(a^{-1})$。
\Needspace{9\baselineskip}
\item\label{item:ch02-laplace-bessel} 对 $a>1$，$\int_0^\infty e^{-ax}J_0(x)\,\mathrm{d}x=(s^2+1)^{-1/2}$，其中
\begin{definition*}[零阶 Bessel 函数]
$\displaystyle J_0(x)=\sum_0^\infty\frac{(-1)^n x^{2n}}{4^n(n!)^2}$ 为零阶 Bessel 函数。
\end{definition*}
\end{enumerate}
\end{exercise}
\end{document}
